\chapter{2012年四川卷（文）}

\section{单选题}

\begin{problem}
设集合 \(A = \{a, b\}\)，\(B = \{b, c, d\}\)，则 \(A \cup B =\)（\quad）

\choices
    {\(\{b\}\)}
    {\(\{b, c, d\}\)}
    {\(\{a, c, d\}\)}
    {\(\{a, b, c, d\}\)}
\end{problem}

\begin{answer}
D
\end{answer}

\begin{solution}
并集中的元素是两个集合中所有不重复的元素，因此 \(A\cup B=\{a,b,c,d\}\)，故选 D.
\end{solution}

\begin{problem}
\((1 + x)^{7}\) 的展开式中 \(x^{2}\) 的系数是（\quad）

\choices
    {21}
    {28}
    {35}
    {42}
\end{problem}

\begin{answer}
A
\end{answer}

\begin{solution}
二项式展开式中 \(x^2\) 的项为 \(\mathrm C_7^2x^2\)，所以 \(x^2\) 的系数为 \(\mathrm C_7^2=21\)，故选 A.
\end{solution}

\begin{problem}
交通管理部门为了解机动车驾驶员（简称驾驶员）对某新法规的知晓情况，对甲，乙，丙，丁四个社区做分层抽样调查. 假设四个社区驾驶员的总人数为 \(N\)，其中甲社区有驾驶员 \(96\) 人. 若在甲，乙，丙，丁四个社区抽取驾驶员的人数分别为 \(12\)、\(21\)、\(25\)、\(43\)，则这四个社区驾驶员的总人数 \(N\) 为（\quad）

\choices
    {101}
    {808}
    {1212}
    {2012}
\end{problem}

\begin{answer}
B
\end{answer}

\begin{solution}
分层抽样时各层的抽样比相同. 甲社区的抽样比为 \(12/96=1/8\)，所以四个社区的总抽样比也是 \(1/8\). 因此
\[
N=8(12+21+25+43)=8\times101=808
\]
故选 B.
\end{solution}

\begin{problem}
函数 \( y = a^x - a \) （\( a > 0 \)，且 \( a \neq 1 \)） 的图象可能是（\quad）

\choices
    {\choicebitmap[width=0.15\paperwidth]{img/2012/sichuan_liberal/q4_A.png}}
    {\choicebitmap[width=0.15\paperwidth]{img/2012/sichuan_liberal/q4_B.png}}
    {\choicebitmap[width=0.15\paperwidth]{img/2012/sichuan_liberal/q4_C.png}}
    {\choicebitmap[width=0.15\paperwidth]{img/2012/sichuan_liberal/q4_D.png}}
\end{problem}

\begin{answer}
C
\end{answer}

\begin{solution}
函数图象恒过点 \((1,0)\)，且 \(y(0)=1-a\). 当 \(a>1\) 时，函数单调递增，且 \(y(0)<0\)；当 \(0<a<1\) 时，函数单调递减，且 \(y(0)>0\)，并且当 \(x\to+\infty\) 时 \(y\to-a<0\). 对照四个图象，只有 C 同时满足恒过 \((1,0)\)、与 \(y\) 轴交点在 \(x\) 轴上方且递减，故选 C.
\end{solution}

\begin{problem}
如图，正方形 \(ABCD\) 的边长为 \(1\)，延长 \(BA\) 至 \(E\)，使 \(AE = 1\)，连接 \(EC\)，\(ED\)，则 \(\sin \angle CED =\)（\quad）

\bitmapfigure[max width=0.36\linewidth]{img/2012/sichuan_liberal/q5_fig1.png}

\choices
    {\(\frac{3\sqrt{10}}{10}\)}
    {\(\frac{\sqrt{10}}{10}\)}
    {\(\frac{\sqrt{5}}{10}\)}
    {\(\frac{\sqrt{5}}{15}\)}
\end{problem}

\begin{answer}
B
\end{answer}

\begin{solution}
取正方形边长为 \(1\)，建立平面直角坐标系，令 \(A(0,0)\)，\(B(1,0)\)，\(C(1,1)\)，\(D(0,1)\)，则 \(E(-1,0)\). 于是
\(EC=\sqrt{2^2+1^2}=\sqrt5\)，\(ED=\sqrt{1^2+1^2}=\sqrt2\)，\(CD=1\).
由余弦定理，
\[
\cos\angle CED=\frac{EC^2+ED^2-CD^2}{2EC\cdot ED}
=\frac{5+2-1}{2\sqrt5\sqrt2}=\frac3{\sqrt{10}}
\]
因此
\[
\sin\angle CED=\sqrt{1-\cos^2\angle CED}
=\sqrt{1-\frac9{10}}=\frac{\sqrt{10}}{10}
\]
故选 B.
\end{solution}

\begin{problem}
下列命题正确的是（\quad）

\choices
    {若两条直线和同一个平面所成的角相等，则这两条直线平行}
    {若一个平面内有三个点到另一个平面的距离相等，则这两个平面平行}
    {若一条直线平行于两个相交平面，则这条直线与这两个平面的交线平行}
    {若两个平面都垂直于第三个平面，则这两个平面平行}
\end{problem}

\begin{answer}
C
\end{answer}

\begin{solution}
选项 A 不成立. 例如在平面 \(z=0\) 内取 \(x\) 轴和 \(y\) 轴，两条直线与该平面所成的角都为 \(0\)，但它们不平行.

选项 B 不成立. 取两个平面为 \(\pi_1:z=0\) 与 \(\pi_2:y=1\)，在 \(\pi_1\) 内取三点 \((0,0,0)\)、\((1,0,0)\)、\((2,0,0)\). 三点到 \(\pi_2\) 的距离都为 \(1\)，而 \(\pi_1\) 与 \(\pi_2\) 相交，故不平行.

选项 D 不成立. 平面 \(x=0\) 与平面 \(y=0\) 都垂直于平面 \(z=0\)，但前两个平面相交于 \(z\) 轴，故不平行.

对于选项 C，设直线的方向向量为 \(\bs v\). 直线平行于两个相交平面，说明 \(\bs v\) 同时属于这两个平面的方向空间；两个方向空间的交集正是两平面交线的方向空间. 因而该直线与两平面的交线平行. 故选 C.
\end{solution}

\begin{problem}
设 \(\bs{a}\)，\(\bs{b}\) 都是非零向量. 下列四个条件中，使 \( \frac{\bs{a}}{|\bs{a}|} = \frac{\bs{b}}{|\bs{b}|} \) 成立的充分条件是（\quad）

\choices
    {\(|\bs{a}| = |\bs{b}|\) 且 \(\bs{a}\parallel\bs{b}\)}
    {\(\bs{a} = -\bs{b}\)}
    {\(\bs{a}\parallel\bs{b}\)}
    {\(\bs{a} = 2\bs{b}\)}
\end{problem}

\begin{answer}
D
\end{answer}

\begin{solution}
由于 \(\bs a\)，\(\bs b\) 均为非零向量，若 \(\bs a=2\bs b\)，则 \(|\bs a|=2|\bs b|\)，从而
\[
\frac{\bs a}{|\bs a|}=\frac{2\bs b}{2|\bs b|}=\frac{\bs b}{|\bs b|}
\]
选项 A 中平行向量可能反向，选项 B 中两向量反向，选项 C 中平行向量也可能反向，均不能保证单位向量相等. 故选 D.
\end{solution}

\begin{problem}
若变量 \(x\)，\(y\) 满足约束条件 \(\begin{cases}
x - y \geqslant -3,\\
x + 2y \leqslant 12,\\
2x + y \leqslant 12,\\
x \geqslant 0,\\
y \geqslant 0,
\end{cases}\) 则 \(z = 3x + 4y\) 的最大值是（\quad）

\choices
    {12}
    {26}
    {28}
    {33}
\end{problem}

\begin{answer}
C
\end{answer}

\begin{solution}
由 \(x\geq0\)、\(y\geq0\)，可在第一象限内考察边界. 直线 \(x-y=-3\) 与 \(x=0\) 交于 \((0,3)\)，与 \(x+2y=12\) 交于 \((2,5)\). 直线 \(x+2y=12\) 与 \(2x+y=12\) 交于 \((4,4)\)，而 \(2x+y=12\) 与 \(y=0\) 交于 \((6,0)\). 再加上原点，约束区域的顶点为
\((0,0)\)，\((0,3)\)，\((2,5)\)，\((4,4)\)，\((6,0)\).

线性目标函数在有界多边形上的最大值必在顶点取得. 逐点代入 \(z=3x+4y\)，得到
\(z(0,0)=0\)，\(z(0,3)=12\)，\(z(2,5)=26\)，\(z(4,4)=28\)，\(z(6,0)=18\).
所以 \(z_{\max}=28\)，故选 C.
\end{solution}

\begin{problem}
已知抛物线关于 \(x\) 轴对称，它的顶点在坐标原点 \(O\)，并且经过点 \(M(2, y_0)\). 若点 \(M\) 到该抛物线焦点的距离为 3，则 \(|OM| =\)（\quad）

\choices
    {\(2\sqrt{2}\)}
    {\(2\sqrt{3}\)}
    {4}
    {\(2\sqrt{5}\)}
\end{problem}

\begin{answer}
B
\end{answer}

\begin{solution}
设抛物线方程为 \(y^2=4px\)，则焦点为 \(F(p,0)\)，且非退化性给出 \(p\ne0\). 因为 \(M(2,y_0)\) 在抛物线上，所以 \(y_0^2=8p\)，从而 \(p>0\). 由 \(MF=3\) 得
\[
9=(2-p)^2+y_0^2=(2-p)^2+8p=(2+p)^2
\]
又因 \(p>0\)，所以 \(2+p>0\)，从而 \(2+p=3\)，即 \(p=1\). 于是
\[
|OM|=\sqrt{2^2+y_0^2}=\sqrt{12}=2\sqrt3
\]
故选 B.
\end{solution}

\begin{problem}
如图，半径为 \(R\) 的半球 \(O\) 的底面圆 \(O\) 在平面 \(\alpha\) 内，过点 \(O\) 作平面 \(\alpha\) 的垂线交半球面于点 \(A\)，过点 \(O\) 的直径 \(CD\) 与平面 \(\alpha\) 成 \(45^{\circ}\) 的平面与半球面相交，所得交线上到平面 \(\alpha\) 的距离最大的点为 \(B\)，该交线上的一点 \(P\) 满足 \(\angle BOP=60^{\circ}\)，则 \(A\)，\(P\) 两点间的球面距离为（\quad）

\bitmapfigure[max width=0.92\linewidth]{img/2012/sichuan_liberal/q10_fig1.png}

\choices
    {\(R\arccos\frac{\sqrt{2}}{4}\)}
    {\(\frac{\pi R}{4}\)}
    {\(R\arccos\frac{\sqrt{3}}{3}\)}
    {\(\frac{\pi R}{3}\)}
\end{problem}

\begin{answer}
A
\end{answer}

\begin{solution}
设平面 \(\alpha\) 为水平面，过球心 \(O\) 且与 \(\alpha\) 成 \(45^{\circ}\) 的截面为平面 \(\beta\). 在该截圆中，点 \(B\) 到平面 \(\alpha\) 的距离最大，所以 \(OB\) 在平面 \(\beta\) 内垂直于 \(\alpha\cap\beta\)，从而 \(OB\) 与平面 \(\alpha\) 成 \(45^{\circ}\).

取 \(\alpha\cap\beta\) 的方向为水平坐标轴，则半径向量 \(OB\) 的竖直分量为 \(R\sin45^{\circ}\). 由于 \(P\) 在同一截圆上且 \(\angle BOP=60^{\circ}\)，向量 \(OP\) 在该竖直方向上的分量为
\[
R\sin45^{\circ}\cos60^{\circ}=\frac{\sqrt2}{4}R
\]
而 \(OA\) 垂直于 \(\alpha\)，故
\[
\cos\angle AOP=\frac{\overrightarrow{OA}\cdot\overrightarrow{OP}}{R^2}=\frac{\sqrt2}{4}
\]
球面上两点间的距离等于半径乘以对应的圆心角，所以 \(A\)，\(P\) 两点间的球面距离为
\[
R\angle AOP=R\arccos\frac{\sqrt2}{4}
\]
故选 A.
\end{solution}

\begin{problem}
方程 \(ay = b^2 x^2 + c\) 中的 \(a\)，\(b\)，\(c \in \{-2, 0, 1, 2, 3\}\)，且 \(a\)，\(b\)，\(c\) 互不相同. 在所有这些方程所表示的曲线中，不同的抛物线共有（\quad）

\choices
    {28 条}
    {32 条}
    {36 条}
    {48 条}
\end{problem}

\begin{answer}
B
\end{answer}

\begin{solution}
要表示抛物线，必须有 \(a\ne0\) 且 \(b\ne0\)，此时方程可写为
\[
y=\frac{b^2}{a}x^2+\frac ca
\]
其中 \(a\) 可取 \(-2\)，\(1\)，\(2\)，\(3\)，而 \(b^2\) 只有 \(1\)，\(4\)，\(9\) 三种值，分别对应 \(b=1\)、\(b=\pm2\)、\(b=3\). 固定 \(a\) 和 \(b^2\) 后，允许的不同 \(c\) 产生不同的抛物线；当 \(b^2=4\) 时，\(b=2\) 与 \(b=-2\) 只改变 \(b\) 的符号，产生同一个方程，因此应合并计数. 下面各组的二次项系数 \(b^2/a\) 彼此不同，所以不同 \(a\) 或 \(b^2\) 的组之间不会重复.

当 \(a=-2\) 时，\(b^2=1\) 对应 \(b=1\)，可取 \(c\in\{0,2,3\}\)，有 \(3\) 个；\(b^2=4\) 时 \(b=-2\) 与 \(a\) 相同，只能取 \(b=2\)，可取 \(c\in\{0,1,3\}\)，有 \(3\) 个；\(b^2=9\) 对应 \(b=3\)，可取 \(c\in\{0,1,2\}\)，有 \(3\) 个，共 \(9\) 条.

当 \(a=1\) 时，\(b^2=1\) 不可取；\(b^2=4\) 时，\(b=2\) 或 \(b=-2\)，合并后可取 \(c\in\{0,-2,2,3\}\)，有 \(4\) 个；\(b^2=9\) 时可取 \(c\in\{0,-2,2\}\)，有 \(3\) 个，共 \(7\) 条.

当 \(a=2\) 时，\(b^2=1\) 可取 \(c\in\{-2,0,3\}\)，有 \(3\) 个；\(b^2=4\) 时 \(b=2\) 不可取，只能取 \(b=-2\)，还需排除 \(a=2\) 和 \(b=-2\)，故 \(c\in\{0,1,3\}\)，有 \(3\) 个；\(b^2=9\) 时可取 \(c\in\{-2,0,1\}\)，有 \(3\) 个，共 \(9\) 条.

当 \(a=3\) 时，\(b^2=1\) 可取 \(c\in\{-2,0,2\}\)，有 \(3\) 个；\(b^2=4\) 时 \(b=2\) 或 \(b=-2\)，合并后可取 \(c\in\{-2,0,1,2\}\)，有 \(4\) 个；\(b^2=9\) 不可取，共 \(7\) 条.

因此不同的抛物线共有
\[
9+7+9+7=32\text{ 条}
\]
故选 B.
\end{solution}

\begin{problem}
设函数 \( f(x) = (x - 3)^3 + x - 1 \)，\( \{a_n\} \) 是公差不为 0 的等差数列，\( f(a_1) + f(a_2) + \dots + f(a_7) = 14 \)，则 \( a_1 + a_2 + \dots + a_7 = \)（\quad）

\choices
    {0}
    {7}
    {14}
    {21}
\end{problem}

\begin{answer}
D
\end{answer}

\begin{solution}
设等差数列的公差为 \(d\ne0\)，并令 \(a_4=3+t\).
\(a_{4+k}=3+t+kd\)，\((k=-3,-2,\ldots,3)\).
由 \(f(x)=(x-3)^3+x-1\)，有
\[
\begin{aligned}
\sum_{k=-3}^{3}f(a_{4+k})
&=\sum_{k=-3}^{3}(t+kd)^3+\sum_{k=-3}^{3}(3+t+kd)-7\\
&=7t^3+3td^2\sum_{k=-3}^{3}k^2+7t+14\\
&=7t^3+84td^2+7t+14.
\end{aligned}
\]
这里利用了 \(\sum_{k=-3}^{3}k=\sum_{k=-3}^{3}k^3=0\)，以及
\(\sum_{k=-3}^{3}k^2=28\).
题设给出该和为 \(14\)，故
\[
7t(t^2+12d^2+1)=0
\]
由于 \(t^2+12d^2+1>0\)，所以 \(t=0\)，即 \(a_4=3\). 于是
\[
a_1+a_2+\cdots+a_7=7a_4=21
\]
故选 D.
\end{solution}

\section{填空题}

\begin{problem}
函数 \( f(x) = \frac{1}{\sqrt{1 - 2x}} \) 的定义域是\fillinblank{}（用区间表示）.
\end{problem}

\begin{answer}
\(\left(-\infty, \frac{1}{2}\right)\)
\end{answer}

\begin{solution}
由于根式在分母中，必须满足
\[
1-2x>0
\]
解得 \(x<\frac12\)，所以函数的定义域为 \(\left(-\infty,\frac12\right)\).
\end{solution}

\begin{problem}
如图，在正方体 \(ABCD - A_{1}B_{1}C_{1}D_{1}\) 中，\(M\)，\(N\) 分别是 \(CD\)，\(CC_{1}\) 的中点，则异面直线 \(A_{1}M\) 与 \(DN\) 所成角的大小是\fillinblank{}.
\bitmapfigure[max width=0.36\linewidth]{img/2012/sichuan_liberal/q14_fig1.png}
\end{problem}

\begin{answer}
\(90^{\circ}\)
\end{answer}

\begin{solution}
设正方体棱长为 \(1\)，建立直角坐标系，取
\(A(0,0,0)\)，\(B(1,0,0)\)，\(C(1,1,0)\)，\(D(0,1,0)\)，\(A_1(0,0,1)\).
由 \(M\)，\(N\) 分别为 \(CD\)，\(CC_1\) 的中点，得 \(M\left(\frac12,1,0\right)\)，\(N\left(1,1,\frac12\right)\).
于是可取两条异面直线的方向向量为 \(\overrightarrow{A_1M}=\left(\frac12,1,-1\right)\)，\(\overrightarrow{DN}=\left(1,0,\frac12\right)\).
它们的数量积为
\[
\overrightarrow{A_1M}\cdot\overrightarrow{DN}=\frac12-\frac12=0
\]
所以两直线所成的角为 \(90^{\circ}\).
\end{solution}

\begin{problem}
椭圆 \(\frac{x^2}{a^2} + \frac{y^2}{5} = 1\)（\(a\) 为定值，且 \(a > \sqrt{5}\)）的左焦点为 \(F\)，直线 \(x = m\) 与椭圆相交于点 \(A\)，\(B\)，\(\triangle FAB\) 的周长的最大值是 12，则该椭圆的离心率是\fillinblank{}.
\end{problem}

\begin{answer}
\(\frac{2}{3}\)
\end{answer}

\begin{solution}
记椭圆的半长轴、半短轴和半焦距分别为 \(a\)，\(b=\sqrt5\)，\(c=\sqrt{a^2-5}\). 设直线 \(x=m\) 与椭圆交于 \(A(m,y)\)，\(B(m,-y)\)，其中 \(y=b\sqrt{1-\frac{m^2}{a^2}}\). 由于 \(F(-c,0)\)，有
\[
FA=FB=\sqrt{(m+c)^2+y^2}
=\sqrt{a^2+2mc+\frac{c^2m^2}{a^2}}
=a+\frac{cm}{a}
\]
因此三角形 \(FAB\) 的周长为
\[
L=2\left(a+\frac{cm}{a}\right)+2b\sqrt{1-\frac{m^2}{a^2}}
\]
令 \(t=\frac ma\)，则 \(-1\leq t\leq1\)，从而
\[
L=2\left(a+ct+b\sqrt{1-t^2}\right)
\]
由柯西不等式，
\[
ct+b\sqrt{1-t^2}
\leq\sqrt{c^2+b^2}\sqrt{t^2+1-t^2}=a
\]
当 \(t=c/a\) 时等号成立，且 \(0<c/a<1\)，该取值在 \([-1,1]\) 内. 因此周长最大值为 \(4a\). 由 \(4a=12\)，得 \(a=3\)，所以
\[
e=\frac ca=\frac{\sqrt{a^2-5}}a=\frac{\sqrt4}{3}=\frac23
\]
\end{solution}

\begin{problem}
设 \(a\)，\(b\) 为正实数. 现有下列命题：
\begin{circlelist}
\item 若 \(a^2 - b^2 = 1\)，则 \(a - b < 1\)；
\item 若 \(\frac{1}{b} -\frac{1}{a} = 1\)，则 \(a - b < 1\)；
\item 若 \(|\sqrt{a} - \sqrt{b}| = 1\)，则 \(|a - b| < 1\)；
\item 若 \( |a^3 - b^3| = 1 \)，
\end{circlelist}
则 \( |a - b| < 1 \). 其中的真命题有\fillinblank{}（写出所有真命题的编号）.
\end{problem}

\begin{answer}
\(\circled{1}\circled{4}\)
\end{answer}

\begin{solution}
对于命题 \(\circled{1}\)，由 \(a^2-b^2=1\) 得
\[
(a-b)(a+b)=1
\]
又因 \(a\)，\(b>0\)，有 \(a^2=1+b^2>1\)，所以 \(a>1\)，并且 \(a>b\). 因此 \(a+b>1\)，从而 \(a-b=1/(a+b)<1\)，命题 \(\circled{1}\) 正确.

命题 \(\circled{2}\) 中，取 \(a=2\)，\(b=\frac23\)，则
\(\frac1b-\frac1a=\frac32-\frac12=1\)，\(a-b=\frac43>1\)
所以命题 \(\circled{2}\) 错误.

命题 \(\circled{3}\) 中，取 \(a=4\)，\(b=1\)，则 \(|\sqrt a-\sqrt b|=1\)，但 \(|a-b|=3>1\)，所以命题 \(\circled{3}\) 错误.

对于命题 \(\circled{4}\)，不妨设 \(a>b>0\). 由
\[
a^3-b^3=(a-b)(a^2+ab+b^2)=1
\]
得
\[
a-b=\frac1{a^2+ab+b^2}
\]
又因为 \(a^3=b^3+1>1\)，所以 \(a>1\)，从而 \(a^2+ab+b^2>a^2>1\)，故
\(a-b<1\). 若 \(b>a>0\)，交换两者可得
\(b^3-a^3=1\)，从而
\(b-a=1/(b^2+ab+a^2)<1\). 因此无论哪一个数较大，都有
\(|a-b|<1\)，命题 \(\circled{4}\) 正确. 真命题的编号为 \(\circled{1}\)，\(\circled{4}\).
\end{solution}

\section{解答题}

\begin{problem}
某居民小区有两个相互独立的安全防范系统（简称系统）\(A\) 和 \(B\)，系统 \(A\) 和 \(B\) 在任意时刻发生故障的概率分别为 \( \frac{1}{10} \) 和 \(p\).
\begin{enumerate}
\item 若在任意时刻至少有一个系统不发生故障的概率为 \(\frac{49}{50}\)，求 \(p\) 的值；
\item 求系统 \(A\) 在 3 次相互独立的检测中不发生故障的次数大于发生故障的次数的概率.
\end{enumerate}
\end{problem}

\begin{solution}
\begin{enumerate}
\item 设系统 \(A\)，\(B\) 发生故障的事件分别为 \(A_0\)，\(B_0\). 由题意，
\(P(A_0)=\frac1{10}\)，\(P(B_0)=p\).
两系统相互独立，且至少有一个系统不发生故障的对立事件是两个系统同时发生故障，所以
\[
1-P(A_0\cap B_0)=\frac{49}{50}
\]
因此
\(1-\frac1{10}p=\frac{49}{50}\)，\(\frac1{10}p=\frac1{50}\)，\(p=\frac15\).
\item 设系统 \(A\) 在 3 次检测中不发生故障的次数为 \(X\)，则
\[
X\sim B\left(3,\frac9{10}\right)
\]
不发生故障的次数大于发生故障的次数，等价于 \(X>3-X\)，即 \(X=2\) 或 \(X=3\). 所求概率为
\[
\begin{aligned}
P(X=2)+P(X=3)
&=\mathrm C_3^2\left(\frac9{10}\right)^2\frac1{10}+\left(\frac9{10}\right)^3\\
&=\frac{243}{250}.
\end{aligned}
\]
\end{enumerate}
\end{solution}

\begin{problem}
已知函数 \( f(x) = \cos^2\frac{x}{2} - \sin \frac{x}{2}\cos \frac{x}{2} - \frac{1}{2} \).
\begin{enumerate}
\item 求函数 \( f(x) \) 的最小正周期和值域.
\item 若 \( f(\alpha)=\frac{3\sqrt{2}}{10} \)，求 \( \sin 2\alpha \) 的值.
\end{enumerate}
\end{problem}

\begin{solution}
\begin{enumerate}
\item 利用二倍角公式，
\[
f(x)=\frac{1+\cos x}{2}-\frac{\sin x}{2}-\frac12
=\frac12(\cos x-\sin x)
=\frac{\sqrt2}{2}\cos\left(x+\frac\pi4\right)
\]
因此 \(f(x)\) 的最小正周期为 \(2\pi\)，值域为 \(\left[-\frac{\sqrt2}{2},\frac{\sqrt2}{2}\right]\).
\item 由题设，
\(\frac12(\cos\alpha-\sin\alpha)=\frac{3\sqrt2}{10}\)，\(\cos\alpha-\sin\alpha=\frac{3\sqrt2}{5}\).
两边平方，得
\[
1-\sin2\alpha=\frac{18}{25}
\]
所以 \(\sin2\alpha=\frac7{25}\).
\end{enumerate}
\end{solution}

\begin{problem}
如图，在三棱锥 \(P - ABC\) 中，\(\angle APB = 90^{\circ}\)，\(\angle PAB = 60^{\circ}\)，\(AB = BC = CA\)，点 \(P\) 在平面 \(ABC\) 内的射影 \(O\) 在 \(AB\) 上.
\begin{enumerate}
\item 求直线 \(PC\) 与平面 \(ABC\) 所成角的正弦值；
\item 求二面角 \(B - AP - C\) 的余弦值.
\end{enumerate}
\bitmapfigure[max width=0.42\linewidth]{img/2012/sichuan_liberal/q19_fig1.png}
\end{problem}

\begin{solution}
\begin{enumerate}
\item 取正三角形 \(ABC\) 的边长为 \(1\)，建立空间直角坐标系，令 \(A(0,0,0)\)，\(B(1,0,0)\)，\(C\left(\frac12,\frac{\sqrt3}{2},0\right)\).
因为 \(O\) 在 \(AB\) 上，所以可设 \(O(t,0,0)\)，并设 \(P(t,0,h)\)，其中 \(h>0\). 由 \(\angle PAB=60^{\circ}\)，得 \(\frac{t}{\sqrt{t^2+h^2}}=\cos60^{\circ}=\frac12\)，从而 \(h^2=3t^2\).
又由 \(\angle APB=90^{\circ}\)，有
\(\overrightarrow{PA}\cdot\overrightarrow{PB}=0\)，\(\Longrightarrow\)，\(t(t-1)+h^2=0\)
即 \(h^2=t(1-t)\). 联立得 \(3t^2=t(1-t)\). 由于 \(t>0\)，所以 \(t=\frac14\)，\(h=\frac{\sqrt3}{4}\).

于是 \(\overrightarrow{PC}=\left(\frac14,\frac{\sqrt3}{2},-\frac{\sqrt3}{4}\right)\)，且 \(|PC|=1\).
平面 \(ABC\) 为水平面，所以直线 \(PC\) 与平面 \(ABC\) 所成角的正弦值为方向向量竖直分量的绝对值与向量长度之比，即
\[
\sin\theta=\frac{\sqrt3/4}{1}=\frac{\sqrt3}{4}
\]
\item 平面 \(BAP\) 的方程为 \(y=0\)，可取其法向量 \(\bs{n}_1=(0,1,0)\). 又 \(\overrightarrow{AP}=\left(\frac14,0,\frac{\sqrt3}{4}\right)\)，\(\overrightarrow{AC}=\left(\frac12,\frac{\sqrt3}{2},0\right)\).
因此平面 \(CAP\) 的一个法向量为
\[
\bs{n}_2=\overrightarrow{AP}\times\overrightarrow{AC}
=\frac18(-3,\sqrt3,\sqrt3)
\]
二面角取锐角时，其余弦为两个法向量夹角余弦的绝对值，故
\[
\cos\varphi
=\frac{|\bs{n}_1\cdot\bs{n}_2|}{|\bs{n}_1|\,|\bs{n}_2|}
=\frac{\sqrt3/8}{\sqrt{15}/8}
=\frac{\sqrt5}{5}
\]
\end{enumerate}
\end{solution}

\begin{problem}
已知数列 \(\{a_{n}\}\) 的前 \(n\) 项和为 \(S_{n}\)，常数 \(\lambda > 0\)，且 \(\lambda a_{1}a_{n} = S_{1} + S_{n}\). 对一切正整数 \(n\) 都成立.
\begin{enumerate}
\item 求数列 \( \{a_{n}\} \) 的通项公式；
\item 设 \(a_{1} > 0\)，\(\lambda = 100\). 当 \(n\) 为何值时，数列 \(\left\{\lg \frac{1}{a_n}\right\}\) 的前 \(n\) 项和最大.
\end{enumerate}
\end{problem}

\begin{solution}
\begin{enumerate}
\item 令 \(n=1\)，由题设得
\[
\lambda a_1^2=2a_1
\]
所以 \(a_1=0\) 或 \(a_1=\frac2\lambda\). 若 \(a_1=0\)，则原等式变为 \(S_n=0\). 由 \(a_n=S_n-S_{n-1}\)，其中约定 \(S_0=0\)，可得 \(a_n=0\) 对一切正整数 \(n\) 成立；反过来全零数列确有 \(S_1=S_n=0\)，满足原等式.

若 \(a_1=\frac2\lambda\)，对任意 \(n\geq2\)，分别写出题设中 \(n\) 与 \(n-1\) 时的等式并相减，得
\[
\lambda a_1(a_n-a_{n-1})=a_n
\]
由于 \(\lambda a_1=2\)，所以 \(a_n=2a_{n-1}\)，从而
\(a_n=\frac{2^n}{\lambda}\)，\((n\geq1)\).
还需验证这两个候选数列确实满足原条件：全零数列有 \(S_1=S_n=0\)，故等式两边均为 \(0\)；若 \(a_n=2^n/\lambda\)，则
\(S_n=\frac{2^{n+1}-2}{\lambda}\)，\(\lambda a_1a_n=\frac{2^{n+1}}{\lambda} =\frac2\lambda+\frac{2^{n+1}-2}{\lambda}=S_1+S_n\).
因此数列的通项有两种可能：全零数列，或 \(a_n=\frac{2^n}{\lambda}\).
\item 由 \(a_1>0\) 可知取第二种情形. 令
\[
T_n=\sum_{k=1}^{n}\lg\frac1{a_k}
\]
当 \(\lambda=100\) 时，\(a_k=2^k/100\)，所以
\[
\lg\frac1{a_k}=\lg\frac{100}{2^k}=2-k\lg2
\]
进而
\[
T_n=2n-\frac{n(n+1)}2\lg2
\]
考察相邻两项之差，
\[
T_{n+1}-T_n=2-(n+1)\lg2
\]
因为 \(2^6<100<2^7\)，所以 \(6\lg2<2<7\lg2\). 故当 \(n\leq5\) 时 \(T_{n+1}-T_n>0\)，当 \(n\geq6\) 时 \(T_{n+1}-T_n<0\). 因此 \(T_n\) 在 \(n=6\) 时取得最大值.
\end{enumerate}
\end{solution}

\begin{problem}
如图，动点 \(M\) 与两定点 \(A(-1,0)\)，\(B(1,0)\) 构成 \(\triangle MAB\)，且直线 \(MA\)，\(MB\) 的斜率之积为 \(4\). 设动点 \(M\) 的轨迹为 \(C\).
\begin{enumerate}
\item 求轨迹 \(C\) 的方程；
\item 设直线 \( y = x + m \) （\( m > 0 \)） 与 \( y \) 轴相交于点 \(P\). 与轨迹 \(C\) 相交于点 \(Q\)，\(R\)，且 \( |PQ| < |PR| \)，求 \( \frac{|PR|}{|PQ|} \) 的取值范围.
\end{enumerate}
\bitmapfigure[max width=0.42\linewidth]{img/2012/sichuan_liberal/q21_fig1.png}
\end{problem}

\begin{solution}
\begin{enumerate}
\item 设 \(M(x,y)\). 由于 \(MAB\) 为三角形，\(y\ne0\)，且 \(x\ne\pm1\). 两直线的斜率分别为
\(k_{MA}=\frac{y}{x+1}\)，\(k_{MB}=\frac{y}{x-1}\).
由斜率之积为 \(4\)，得
\(\frac{y^2}{(x+1)(x-1)}=4\)，\(y^2=4(x^2-1)\).
所以轨迹 \(C\) 的方程为
\[
x^2-\frac{y^2}{4}=1
\]
（其中取 \(y\ne0\) 的点；所以轨迹是该双曲线除去两个顶点 \((\pm1,0)\)）.
\item 直线 \(y=x+m\) 与 \(y\) 轴交于 \(P(0,m)\). 将 \(y=x+m\) 代入轨迹方程，得
\((x+m)^2=4(x^2-1)\)，\(3x^2-2mx-(m^2+4)=0\).
记 \(s=\sqrt{m^2+3}\)，则两个交点的横坐标为
\(x_-=\frac{m-2s}{3}<0\)，\(x_+=\frac{m+2s}{3}>0\).
直线上的点 \((x,x+m)\) 到 \(P\) 的距离为 \(\sqrt2|x|\). 由于 \(m>0\)，有 \(x_++x_-=2m/3>0\)，所以 \(x_+>-x_-\)，从而较近的交点对应 \(x_-\)，较远的交点对应 \(x_+\). 因此
\[
\frac{|PR|}{|PQ|}
=\frac{x_+}{-x_-}
=\frac{m+2s}{2s-m}
\]
当 \(m=1\) 时，较近的交点为 \((-1,0)\)，不属于轨迹 \(C\)，故由题设 \(Q\)，\(R\) 均在轨迹 \(C\) 上可知 \(m\ne1\).
令 \(u=\frac{m}{s}\)，则 \(0<u<1\)，并且
\[
\frac{|PR|}{|PQ|}=\frac{u+2}{2-u}
\]
该函数的导数为
\[
\left(\frac{u+2}{2-u}\right)'=\frac4{(2-u)^2}>0
\]
所以它在 \(0<u<1\) 上单调递增. 又
\[
u'(m)=\frac3{(m^2+3)^{3/2}}>0
\]
故 \(m\) 从 \(0\) 增大到 \(+\infty\) 时 \(u\) 严格递增. 当 \(m=1\) 时，\(u=\frac12\)，对应的比值为 \(\frac53\). 排除 \(m=1\) 后，取值范围为
\[
\left(1,\frac53\right)\cup\left(\frac53,3\right)
\]
\end{enumerate}
\end{solution}

\begin{problem}
已知 \(a\) 为正实数，\(n\) 为自然数，抛物线 \(y = -x^{2} + \frac{a^{n}}{2}\) 与 \(x\) 轴正半轴相交于点 \(A\). 设 \(f(n)\) 为该抛物线在点 \(A\) 处的切线在 \(y\) 轴上的截距.
\begin{enumerate}
\item 用 \(a\) 和 \(n\) 表示 \( f(n) \)；
\item 求对所有 \(n\) 都有 \(\frac{f(n) - 1}{f(n) + 1} \geqslant \frac{n}{n + 1}\) 成立的 \(a\) 的最小值；
\item 当 \(0 < a < 1\) 时，比较 \(\frac{1}{f(1) - f(2)} + \frac{1}{f(2) - f(4)} + \cdots + \frac{1}{f(n) - f(2n)}\) 与 \(6 \cdot \frac{f(1) - f(n + 1)}{f(0) - f(1)}\) 的大小，并说明理由.
\end{enumerate}
\end{problem}

\begin{solution}
\begin{enumerate}
\item 抛物线方程为 \(y=-x^2+\frac{a^n}{2}\)，其正半轴交点为
\[
A\left(\frac{a^{n/2}}{\sqrt2},0\right)
\]
在点 \(A\) 处的切线斜率为 \(-2x_A=-\sqrt2a^{n/2}\)，所以切线方程为
\[
y=-2x_A(x-x_A)
\]
令 \(x=0\)，得到 \(y\) 轴截距
\[
f(n)=2x_A^2=a^n
\]
\item 由 \(a>0\)，有 \(f(n)+1=a^n+1>0\)，所以题设不等式等价于
\[
(n+1)(a^n-1)\geq n(a^n+1)
\]
即
\[
a^n\geq2n+1
\]
取 \(n=1\)，必有 \(a\geq3\). 反之，若 \(a=3\)，由二项式不等式
\[
3^n=(1+2)^n\geq1+2n
\]
可知对所有正整数 \(n\) 都满足要求. 因此 \(a\) 的最小值为 \(3\).
\item 当 \(0<a<1\) 时，由 \(f(k)=a^k\)，左边记为 \(L\)，右边记为 \(R\)，则
\[
L=\sum_{k=1}^{n}\frac1{a^k-a^{2k}}
=\sum_{k=1}^{n}\frac1{a^k(1-a^k)}
\]
并且
\[
R=6\frac{a-a^{n+1}}{1-a}=6\sum_{k=1}^{n}a^k
\]
对任意 \(x\in(0,1)\)，令 \(h(x)=x^2(1-x)\)，则
\[
h'(x)=x(2-3x)
\]
所以 \(h\) 在 \(0<x<\frac23\) 上递增，在 \(\frac23<x<1\) 上递减，从而在 \(x=\frac23\) 处取得最大值 \(\frac4{27}\). 因此
\[
6x^2(1-x)\leq\frac89<1
\]
令 \(x=a^k\)，即得
\[
\frac1{a^k(1-a^k)}>6a^k
\]
对 \(k=1\)，\(2\)，\(\ldots\)，\(n\) 求和，得到 \(L>R\)，所以左边的式子大于右边的式子.
\end{enumerate}
\end{solution}
